\paragraph{F.48 (Proved). Recovered Section 53 comparison and the exact cyclotomic rows.}
\textit{Manuscript: Longue Marche, \S53.5A--C, folios 768--779,
cote 140-4 scans 242--262, especially (47), (56), (81)--(94);
Montpellier cote 149, \S1, folios 6--7, scans 13--15, and \S4,
folios 63--65, scans 128--131; cote 147, \S11, scans 143--148.}
The one-dimensional specialization of F.19 on a smooth surface is the sector-filter
functor actually constructed in \S53, not just an abstract equivalence
between two groupoids having fundamental group $\mathbb Z$. Its finite
completion gives the complex comparison of the tangential fibre functor.
The recovered \texttt{CYCLOTOMIC\_KERNEL} register has exact row identifiers
\texttt{B5} and \texttt{C8}. B5 is a slimness statement; C8 is diagonal
boundary nonvanishing together with a counterexample to unrestricted
local injectivity. Their mathematical contents agree with the corrected
statements below, and do not assert a universal converse recognizing
cuspidal sections.

\begin{proof}
\emph{The germ and its coverings.}
Let $x$ be a point of a smooth real surface $X$ and write
$D_x^*$ for the germ of punctured neighbourhoods of $x$.
Let $S_x=(T_xX\setminus0)/\mathbb R_{>0}$ be its circle of positive
half-rays. A covering of a germ means the filtered colimit of the
categories of coverings of its representatives, exactly (47) on scan
242. Small punctured discs form a cofinal system; restriction between
them induces an equivalence on all covering spaces. Therefore
\[
 \operatorname{Rev}(D_x^*)\simeq\mathbb Z\text{-}\mathbf{Set},
 \qquad \operatorname{Rev}_{\mathrm{uni}}(D_x^*)\simeq B\mathbb Z.
 \tag{F48.1}
\]
Here the latter is the groupoid of connected simply connected covering
germs and their isomorphisms. Choosing an orientation identifies its
automorphism group with $\mathbb Z$; intrinsically it is the orientation
module denoted $\overline T$ in (56), scan 248.

The poset $\mathcal A_x$ in (82), scan 256, consists of nonempty
\emph{closed proper connected arcs} of $S_x$, including single points
and excluding the whole circle. Its groupoid completion is
$\Pi_1(S_x)$. Indeed, choose a point in every arc. An inclusion of
arcs determines a unique homotopy class of paths in the larger arc.
Conversely, subdividing any path into finitely many arc neighbourhoods
expresses it by these generators. A homotopy of paths has a finite
subdivision into triangles lying in such neighbourhoods, so the only
relations required are compositions of inclusions and their inverses.
These constructions are inverse up to the evident changes of the
chosen points. This proves (83)--(85), scans 257--258, without assuming
that the arc poset is simply connected.

\emph{Strict thickening and coordinate independence.}
For $A\in\mathcal A_x$, let $C_A$ be its cone in $T_xX$. Choose a
local differentiable chart from the tangent plane to $X$ inducing the
identity on $T_xX$. The germ associated to $A$ is represented by the
filter generated by
\[
  \{v:0<\lVert v\rVert<\epsilon,\ [v]\in A'\},
  \qquad A\subset\operatorname{int}(A'),\quad A'\subsetneq S_x,
 \tag{F48.2}
\]
where $A'$ runs through strict angular enlargements and $\epsilon>0$.
One may take $A'$ open or take interiors of closed strict enlargements;
the generated filters are cofinal. Every such representative is
contractible, by polar coordinates on an angular interval. No
convexity assertion is needed when the interval is wider than a
semicircle.

If $h$ changes the chosen chart and $d h_0=1$, then
$h(v)=v+o(\lVert v\rVert)$ uniformly in direction. Given a strict
enlargement $A'$ choose an intermediate enlargement $A''$ whose closure
is contained in $\operatorname{int}(A')$. For sufficiently small $\delta$,
$h$ sends the sector for $(A'',\delta)$ into the sector for
$(A',\epsilon)$. Apply the same observation to $h^{-1}$.
Thus both charts define the same filter germ. This is precisely the
strict-enlargement construction (89)--(90) on scans 259--260.
Using just a fixed ray would fail: in coordinates $(s,t)$ the
diffeomorphism $(s,t)\mapsto(s,t+s^2)$ has identity derivative at zero
but does not preserve the germ of the positive $s$-axis.

Pulling a covering back to this simply connected sector germ gives a
set of sheets and hence a fibre functor $F_A$ on
$\operatorname{Rev}(D_x^*)$. Inclusions of arcs give canonical
isomorphisms between these fibre functors, by restriction to the
smaller sector germ; their compatibility is the compatibility of
restriction. The associated universal covering germs therefore define
\[
 \Theta:\Pi_1(S_x)\longrightarrow
              \operatorname{Rev}_{\mathrm{uni}}(D_x^*),
 \tag{F48.3}
\]
the functor (81), scan 256. Rotating the direction once lifts, in
angular coordinates, to a translation by one full turn on the
universal covering of a punctured disc. Thus $\Theta$ sends the
generator of $\pi_1(S_x)$ to a generator of the automorphism group
in (F48.1), with the sign fixed by the chosen fibre-functor and
representing-cover convention. Both groupoids are connected and nonempty, so $\Theta$ is
fully faithful and essentially surjective. This proves the
equivalence by the winding criterion in (86bis), scan 258.

\emph{Arcs and tangential points.}
An outgoing differentiable arc with nonzero derivative and tangent
half-ray $\Delta$ eventually lies in every strict sector enlargement
of $\Delta$. It consequently gives a map of germs into (F48.2).
On coverings this identifies its fibre functor with $F_{\{\Delta\}}$,
canonically and compatibly with morphisms of universal covers. These
are (91)--(94), scans 261--262. The normal-direction construction in
F.19 uses these same simply connected sectors and restriction maps;
its rank-one specialization is consequently (F48.3). This verifies
agreement of the actual functors, including their path morphisms.

For a complex algebraic curve, apply the finite covering comparison
and complete the automorphism group to $\widehat{\mathbb Z}(1)$.
In a local parameter $z$, the finite connected covers are cofinally
the Kummer covers $z^{1/n}$; continuation along the sector picks their
compatible sheets. This is the fibre functor in Deligne's
\S15.13--15.27, as identified in F.19. Replacing $z$ by $u z$ changes
the compatible lifts by choices of $n$th roots of the unit $u$;
their ratios are exactly the Kummer transition torsor. Thus the
comparison retains the parameter-change identifications. In
particular it does not invent a canonical map of spaces
$D_x^*\to S_x$, whose absence is explicitly discussed on scan 262.

There is a separate correction on scans 242--243: the speculative
identification with the inverse limit of the \emph{full small topoi}
of all shrinking punctured neighbourhoods is false. Their localic
intersection is empty, so that inverse limit is the degenerate topos,
whereas (F48.1) is not degenerate. The correct objects are the
pro-topos of neighbourhoods, or the category of covering germs.
This is the exact LM8 E.10/E.11 distinction proved in F.12, now
checked against the original wording, including the inserted word
\emph{petits} on scan 243.

\emph{The recovered register B5.}
Write $\mathsf{Slim}(G)$ for the condition
$C_G(H)=1$ for every open $H\subseteq G$. It is preserved by
extensions: the image of a centralizer element is trivial in a slim
quotient, and its remaining element in the kernel centralizes the
open intersection there. It is also preserved by inverse limits of
slim profinite groups with surjective transition maps. A centralizer
element has trivial coordinates, since projections are open and
send the given open subgroup to open subgroups.
The nonabelian finite-rank free profinite case is the case established
in the free-group argument of F.39. More explicitly, an open normal
subgroup has free rank $1+[F:H](\operatorname{rank}F-1)>1$ and trivial
centre. Its centralizer injects into the finite group $F/H$ and is
trivial by torsion-freeness. Rank one is excluded because
$F=\widehat{\mathbb Z}$ is its own centralizer. This proves exactly
the three assertions of the recovered B5, including its correction
to scan 15. The identifier does not designate an inertia-recognition
or cyclotomic-normalization theorem.

Here torsion-freeness can be checked without a separate group-theoretic
black box. If a nontrivial finite-order element $z\in F$ survives in a
finite quotient $F/N$, put $V=N\overline{\langle z\rangle}$. The open
group $V$ is free profinite, so its abelianization is a finite product
of copies of $\widehat{\mathbb Z}$, a torsion-free group. But the
nonzero cyclic quotient $V/N$ generated by $zN$ factors through that
abelianization, while the finite-order element $z$ maps to zero there.
This contradiction proves the assertion used above.

\emph{The recovered register C8.}
Let $U/k$ be a tower of relative hyperbolic curves in characteristic
zero, $U\subset X$ a smooth proper normal-crossings compactification,
and $x$ a rational boundary point with geometric branch set $I$.
The local geometric group is $T^I$, $T=\widehat{\mathbb Z}(1)$,
with action given by cyclotomic scaling and branch permutation
\citep[Proposition~3.4]{CarlsonCorwin2026}. For its map
$j_x:T^I\to\pi_1(U_{\bar k})$, the diagonal $T\to T^I$ has nonzero
image. Here is the additional argument, beyond local normalization.
Hoshi's valuative criterion for a relative hyperbolic curve says that
a map from the fraction field of a strictly henselian trait extends
over the trait precisely when its fundamental-group image is trivial
\citep[Lemma~2.8, pp.~175--176]{Hoshi2014}. Induction up the tower
gives that criterion for $U$: extend the map to the previous stage
first, and then apply the relative-curve criterion at the last stage.

In formal branch coordinates at a geometric point over $x$, take
the arc $z_i=t$ for every boundary parameter, and take the remaining
parameters to be zero. It has generic point in $U$ and closed point
$x$. Pullback of $z_i^{1/n}$ to $t^{1/n}$ shows that its fundamental
group map is exactly the diagonal followed by $j_x$. If this were
zero the criterion would extend the generic-point map to $U$.
Separatedness of $X$ would identify this extension with the original
arc, contradicting its boundary closed point. This proves
nonvanishing. The cyclotomic kernel fixes the diagonal regardless
of branch permutation, and fixes all of $T^I$ when the branches
are rational. The stated cyclotomic centralizer conclusion follows.

Full injectivity fails in the blown-up product already computed in
F.39: at the new crossing the map is
$T^2\to T$, $(a,b)\mapsto a+b$, with kernel generated by $(1,-1)$.
The diagonal maps by multiplication by two and is nonzero. Thus the
stable-model hypotheses in F.39 and the counterexample in C8 are
compatible. A converse boundary-local factorization of all sections
with nonzero cyclotomic centralizer would imply the point-or-cusp
section assertion for curves: apply it to the product of an arbitrary
curve section and a rational cusp section of the triply punctured
line, and project back to the curve. Neither B5 nor C8 supplies this
converse. Carlson and Corwin's Theorem~3.7 and Proposition~4.12
concern closures of the corresponding section sets; dropping those
closure bars would change their statements.
\end{proof}

\emph{Audit evidence.}
The recovered register, the B5 and C8 proof files, the two \S53
transcriptions, and their file hashes are retained in the source-audit
appendix of the package. The decisive original scans 242, 243, 256,
259--262 were inspected visually. Earlier unsuccessful-search notes
are preserved as dated history and superseded here. The separate
\texttt{FLAGGED\_149} register remains the source obligation in F.45.
